%*******************************************************
% Foreword
%*******************************************************
%\renewcommand{\abstractname}{Abstract}
\pdfbookmark[1]{Prefazione}{Prefazione}
% \addcontentsline{toc}{chapter}{\tocEntry{Abstract}}
\begingroup
\let\clearpage\relax
\let\cleardoublepage\relax
\let\cleardoublepage\relax

\chapter*{Prefazione}
Un eserciziario che nasce dal disordinato accumulo di compiti assegnati sugli algoritmi e pensati per sviluppare il pensiero computazionale, da realizzare con diagrammi di flusso. Finalmente dopo quasi cinque anni mi sono deciso a mettere in ordine (almeno) questa parte di tutto il materiale che produco nella mia didattica. Questa raccolta di esercizi mi torna molto utile nell'insegnamento dell'Informatica durante il primo biennio del Liceo Scientifico opzione Scienze Applicate, tuttavia non escludo che possa potenzialmente diventare una risorsa fruibile anche in altri indirizzi dove si vuole imparare il ragionamento algoritmico, eventualmente anche prima di arrivare alle superiori perché, credo, prima si inizia a sviluppare l'abilità del pensiero computazionale, meglio è.

Essendo un eserciziario non ci sarà ombra di teoria nelle pagine che seguono per due motivi: ci sono già tanti libri che spiegano bene e approfonditamente tutte le nozioni necessarie e, in secondo luogo, non trovo strettamente necessario avere un testo per questa parte di teoria, a me piace spiegarla freestyle.

Un ultimo appunto: ci si potrebbe accorgere che non è tutto impostato e scritto per essere rigidamente serio e "professionale", ma è questo il mio stile e il mio modo di rendere più coinvolgente, o almeno non troppo noioso se possibile, questa raccolta di esercizi.

\chapter*{Istruzioni per l'uso}
Ho suddiviso i capitoli di questo eserciziario in tre sezioni:

\begin{description}
	\item[Sezione I] Il primo mattoncino fondamentale: le variabili. E insieme a queste l'interazione fra algoritmo e utente tramite input e output di dati.
	\item[Sezione II] La struttura di controllo della selezione, che sia semplice, annidata o con condizioni più o meno complesse.
	\item[Sezione III] La selezione utilizzata per realizzare l'iterazione negli algoritmi tramite cicli indefiniti e definiti.
\end{description}

Il mio suggerimento è di procedere seguendo l'ordine dei capitoli poiché, almeno nelle mie intenzioni, gli esercizi di ciascun capitolo permettono di acquisire gli strumenti necessari per affrontare gli esercizi del capitolo successivo.

\chapter*{Note didattiche}
Poche\footnote{L'autore ha scritto questa parola prima di proseguire con le spiegazioni annunciate, pertanto con questo termine potrebbe aver espresso una speranza o un'intenzione piuttosto che una verità fattuale.} parole per spiegare alcune scelte didattiche che ho fatto scrivendo questo eserciziaro.

\begin{description}
	\item[Teoria-free] Come già anticipato, essendo un eserciziario non ho inserito teoria. Lo scopo principale di questi esercizi è quello di concentrarsi sulla logica algoritmica per sviluppare il pensiero computazionale, una delle competenze chiave in un percorso di studi scientifici.
	\item[Contesti vari e reali] Non volevo limitarmi ad esercizi astratti, sotto l'aspetto didattico trovo che far incontrare esercizi e realtà sia un ottimo modo per far capire che si ha a che fare con il pensiero computazionale in situazioni quotidiane, aiutando quindi a capire che gli algoritmi possono modellare la realtà per risolvere problemi con cui si ha a che fare in contesti di tutti i giorni.
	\item[Esercizi riciclati] Si incontreranno nel corso dei capitoli esercizi che si ripetono arricchendosi, non è pigrizia, bensì una strategia per far capire come gli algoritmi si possano raffinare e migliorare con gli strumenti che si imparano a padroneggiare andando avanti.
	\item[Piccole sfide] Spesso mi sono trovato in classi dove qualcuno spicca sul pensiero computazionale quindi ho predisposto anche degli esercizi più sfidanti, così nessuno dovrebbe annoiarsi.
\end{description}

\endgroup